\section{The local action and two reflecting arcs}\label{sec:local}
\subsection{The physical return branch}
Place a normal contact pair at $(0,0)$ and $(g,0)$ for the proof only.
In transverse graph coordinates its boundaries are
$(-\psi_0(y),y)$ and $(g+\psi_1(y),y)$, where
$\psi_b(0)=\psi_b'(0)=0$ and $\psi_b''(0)=\kappa_b>0$.
The broken length has a nondegenerate minimum in its intermediate
coordinate: its second derivative there is $2(g^{-1}+\kappa_o)>0$.
The implicit-function theorem gives a unique stationary intermediate
impact for endpoints in a smaller fixed box. Strict channel clearance
identifies this branch with the physical specular trajectory.

Use intrinsic unit-speed boundary parametrizations $x(s)$ and $y(u)$,
with both tangents in the positive transverse direction at zero, and put
\begin{equation}\label{eq:action}
 D(s,u)=|x(s)-y(u)|,\qquad
 W(s,t)=D(s,u(s,t))+D(t,u(s,t)).
\end{equation}
The stationary equation is $D_u(s,u)+D_u(t,u)=0$. On a sufficiently
small rectangle $W$ is smooth, symmetric in its two arguments, and has
positive dispersing twist $-W_{st}>0$. At the normal orbit,
\[
 W(0,0)=2g,\qquad W_{st}(0,0)=-\frac1{2g(1+g\kappa_o)}.
\]
These statements hold uniformly on compact physical families with the
stated geometry and the required fixed derivative bounds.

\begin{lemma}[The residual-time density]\label{lem:density}
On a sufficiently short onset collar the endpoint density is
\begin{equation}\label{eq:density}
 f_T(s,t)=\frac{-W_{st}(s,t)}{2\pi A}(T-W(s,t))_+.
\end{equation}
On a rectangle where $T_1-W>0$, the two density functions at
$T_2>T_1$ recover $W$ and $A$ by \eqref{eq:ratio-main}.
\end{lemma}
\begin{proof}
Let $p$ be the momentum conjugate to source arclength and $r$ the time
before its first impact. Collision-section phase volume is
$\dd s\dd p\dd r$. The generating identity $p=-W_s$ changes this to
$(-W_{st})\dd s\dd t\dd r$. The allowed residual interval is
$0<r<T-W(s,t)$. Choose the collar shorter than a common positive lower
bound for adjacent free flights. Then this interval cannot include a
preceding collision, and the time after the third impact cannot include
a fourth. The strict local minimum and patch isolation keep the active
sublevel in this physical chart. Conversely, every selected near-onset
trajectory has these coordinates. Integrating $r$ and normalizing by
$2\pi A$ proves \eqref{eq:density}.

On a strictly active rectangle set $w=-W_{st}/(2\pi A)>0$.
Then $f_j=w(T_j-W)$ and $f_2-f_1=(T_2-T_1)w$.
Eliminating $w$ proves the action formula, and its recovered mixed
second derivative proves the area formula. Equality of measures
uniquely determines their continuous density representatives.
\end{proof}

\begin{proposition}[Exact local information equivalence]\label{prop:lens}
On the physical class and a fixed strictly active rectangle, the two
endpoint density functions at known absolute times are equivalent to
$(W,A)$. In particular they determine the local canonical scattering
relation
\begin{equation}\label{eq:lens}
       (s,-W_s(s,t))\longmapsto(t,W_t(s,t)).
\end{equation}
Conversely, that relation, one absolute action value, and $A$ determine
the densities on the rectangle. No extension to a global exterior
travelling-time relation is asserted.
\end{proposition}
\begin{proof}
The forward and inverse formulas are Lemma~\ref{lem:density}.
The positive twist makes $t\mapsto-W_s(s,t)$ locally invertible.
Thus \eqref{eq:lens} is a canonical graph; its generating one-form is
$-p_-\dd s+p_+\dd t=\dd W$. On the connected rectangle it determines
$W$ up to a constant, fixed by one value. Substitution in
\eqref{eq:density} gives the converse. Here the given relation is assumed
physical, so its generating one-form is exact. Arbitrary positive
functions satisfying the quotient formula are not being declared
physical endpoint laws.
\end{proof}

This makes explicit both the richness and the origin of the data.
First derivatives of travel time recover endpoint directions in
scattering geometry; compare Noakes--Stoyanov~\cite[Lemma 4]{NS} and
Gurfinkel--Noakes--Stoyanov~\cite[Lemma 7]{GNS}. The affine elimination
is not a new abstract lens-rigidity principle. What remains to prove
here is an explicit inverse for two unknown reflecting arcs in their
intrinsic coordinates and its periodic descent with volume calibration.

\subsection{The curvature identities}
On the diagonal write $u(s)=u(s,s)$ and define
\begin{equation}\label{eq:diagonal}
 \ell(s)=\tfrac12W(s,s),\qquad a(s)=W_s(s,s),\qquad
 v(s)=\sqrt{1-a(s)^2}.
\end{equation}
Shrink the interval so that $v>0$. Let $T=x'$ and $N$ be the source
unit tangent and inward normal, with $T'=kN$ and $N'=-kT$.
The vector $e=(x-y)/\ell$ then satisfies $e=aT+vN$.

\begin{lemma}[Diagonal Schur complement]\label{lem:curvature}
The source curvature and the opposite curvature at the nearest point are
\begin{align}
 k(s)&=\frac{W_{ss}(s,s)-W_{st}(s,s)-v(s)^2/\ell(s)}{v(s)},
                                                    \label{eq:curvature-source}\\
 k_o(u(s))&=\frac1{\ell(s)}
       \left(\frac{v(s)^2}{-2\ell(s)W_{st}(s,s)}-1\right),
                                                    \label{eq:curvature-other}\\
 u'(s)&=\frac{v(s)}{1+\ell(s)k_o(u(s))}>0.          \label{eq:foot-speed}
\end{align}
\end{lemma}
\begin{proof}
At $s=t$ the stationary equation is $D_u=0$, selecting the nearest
point on the opposite strictly convex patch. Orient its tangent $T_o$
continuously so that $T\cdot T_o=v>0$. Direct differentiation gives
\[
 D_{ss}=v^2/\ell+kv,\qquad D_{su}=-v/\ell,\qquad
 D_{uu}=1/\ell+k_o.
\]
For the last identity, $y''=-k_oe$ at the nearest point; this is the
external, dispersing sign. The stationary elimination of the intermediate
variable gives
\begin{equation}\label{eq:schur}
 W_{ss}=D_{ss}-\frac{D_{su}^2}{2D_{uu}},\qquad
 W_{st}=-\frac{D_{su}^2}{2D_{uu}}
        =-\frac{v^2}{2\ell(1+\ell k_o)}.
\end{equation}
Subtracting the first two expressions gives
\eqref{eq:curvature-source}; solving the last gives
\eqref{eq:curvature-other}. Differentiation of $D_u(s,u(s))=0$
yields \eqref{eq:foot-speed}.
\end{proof}

\begin{proof}[Proof of Theorem~\ref{thm:local-main}, exact assertion]
Recover $W$ and $A$ from Lemma~\ref{lem:density}. Choose
$x(0)=0$, $T(0)=(0,1)$ and $N(0)=(-1,0)$ and solve the Frenet
system with the recovered $k$. This determines the source arc. The
opposite arc in the same plane is
\begin{equation}\label{eq:foot}
 y(u(s))=x(s)-\ell(s)\{a(s)T(s)+v(s)N(s)\}.
\end{equation}
By \eqref{eq:foot-speed} it is a regular arc containing
$y(0)=(g,0)$, where $g=W(0,0)/2$. The recovered objects are functions
on intervals, not merely their Taylor series. Simultaneous arclength
reversal reflects the entire pair, not its members separately. The
remaining choices are one translation and orthogonal transformation.
\end{proof}

\begin{proposition}[Conditional local Lipschitz inverse]\label{prop:local-stable}
Fix nested endpoint rectangles $Q_0\Subset Q$, a diagonal subinterval,
and positive constants $M,g_*,A_*,v_*,w_*,\tau_*,m_*$. Consider smooth
physical pairs for which, in their normal contact frames, both intrinsic
boundary parametrizations have $C^3$ norm at most $M$ on fixed slightly
larger arcs, $g_*,A_*\le g,A\le M$, $|T_j|\le M$,
$T_2-T_1\ge\tau_*$, $T_1-W\ge m_*$ on $Q$,
$-W_{st}\ge w_*$ on $Q$, $v\ge v_*$ on the diagonal, and
$\|f_j\|_{C^2(Q)}\le M$. The patches, isolation and adjacent free-flight
margins are uniform, as is positive curvature. For two such pairs at
the same times put
\[
        e=\max_{j=1,2}\|f_j-\widetilde f_j\|_{C^2(Q)}.
\]
After a permitted common reflection, their gaps, areas, and intrinsic
$C^2$ boundary arcs on smaller fixed intervals differ by at most $Ce$.
The constant depends on all the stated bounds and margins, not on a
contact-jet order. This is an estimate between physical density pairs,
not an inverse defined on every pair of $C^2$ functions.
\end{proposition}
\begin{proof}
The density difference is bounded below by
$\tau_*w_* /(2\pi M)$. Products and reciprocals are Lipschitz on the
corresponding bounded $C^2$ sets. Formula~\eqref{eq:ratio-main} gives
$C^2$ action error $Ce$; its area expression gives scalar error $Ce$.
The diagonal formulas give $Ce$ source-curvature error and opposite
curvature error as functions of $s$, since all their denominators have
positive margins. The Frenet integral equations and Gronwall's inequality
give $Ce$ error for the source arc through its second derivative.

For the opposite arc, differentiating \eqref{eq:foot} twice would
unnecessarily require a third derivative of $W$. Instead integrate
\eqref{eq:foot-speed}. Its speed has positive upper and lower bounds,
and its error is $Ce$. Consequently $u(s)$ and its inverse have error
$Ce$ on smaller intervals. The physical $C^3$ prior bounds the derivative
of opposite curvature in intrinsic arclength. Thus the two opposite
curvature functions differ by $Ce$ at the same $u$. Their Frenet systems,
with the recovered gap and common initial tangent, give the required
$C^2$ estimate. This uses exactly the physical prior stated in the
proposition and no derivative of an arbitrary noisy curvature estimate.
\end{proof}
